%========================================================================
\section{Construction --- Un pipeline de Skills}
% ⏱ ~4 min  |  Le « ce que j'ai construit ». Montrer l'ARCHITECTURE, pas le code.
%========================================================================

\begin{frame}{Automatiser la production de la documentation}
  % 🎤 ~2 min. Documenter 6 pages à la main = plusieurs semaines + incohérences.
  % Solution : une chaîne de 4 Skills Claude, calibrée sur l'étalon avant/après.
  \placeholder{3.0cm}{figure 2.1 du rapport : pipeline site-mapper $\to$ to-tnr $\to$ to-confluence / to-prompt}
  \vskip 4pt
  \footnotesize
  \begin{itemize}
    \item \texttt{site-mapper} --- cartographie exhaustive d'une page (via Playwright)
    \item \texttt{site-mapper-to-tnr} --- rationalise en fonctionnalités TNR au bon niveau
    \item \texttt{tnr-to-confluence} --- export Confluence (DOCX + CSV d'index)
    \item \texttt{tnr-to-prompt} --- transforme la doc en prompt exécutable par un agent
  \end{itemize}
\end{frame}

\begin{frame}{Deux partis pris d'ingénieur}
  % 🎤 ~2 min. Ce sont MES décisions d'architecture (pas l'IA). Bien les assumer.
  \begin{columns}[T]
    \column{.5\textwidth}
    \begin{block}{Séparation en 4 Skills}
      \footnotesize
      Responsabilité unique par étape $\Rightarrow$ on \alert{relance l'une sans les autres}.
      Générique : réapplicable à d'autres sites Shopify/SFCC en n'adaptant que l'entrée.
    \end{block}
    \column{.5\textwidth}
    \begin{alertblock}{Rétro-spec assumée}
      \footnotesize
      Faute de spec, on reconstruit depuis la prod.
      Limite connue : un bug en prod peut devenir une norme.
      Conséquence : tout écart pré-prod / prod sur un élément non modifié = bug.
    \end{alertblock}
  \end{columns}
  \vskip 6pt
  % 🎤 Maintenance : matrice par page dans Confluence, statut « Impacté MEP » avant chaque livraison.
\end{frame}
